\section*{Hoofdstuk \ref{ch:g7:short-circuits-safety} --- Kortsluiting en elektrische veiligheid}

\begin{solution}{exo:g7:short-circuits-safety:1}
Tussen twee wegen die dezelfde punten verbinden, dringt de stroom zich
op de weg die hem het minst tegenwerkt. Een werkende lamp werkt de
mars tegen; een blote dikke draad nauwelijks --- dus laat de mars de
lamp in de steek voor de draad.
\end{solution}

\begin{solution}{exo:g7:short-circuits-safety:2}
\omterm{def:g7:short-circuits-safety:device}{Kortsluiting} van een apparaat: een omleiding over één apparaat --- dat
stopt met werken terwijl de lus doorgaat. \omterm{def:g7:short-circuits-safety:device}{Kortsluiting} van de
\omterm{def:g3:first-electric-circuit:battery}{batterij}: een blote weg die de eigen polen van de
\omterm{def:g3:first-electric-circuit:battery}{batterij} verbindt zonder enig apparaat --- het ernstigste geval: niets
werkt de mars tegen, en de stormloop verhit draad en
\omterm{def:g3:first-electric-circuit:battery}{batterij} tot verwoesting en brand.
\end{solution}

\begin{solution}{exo:g7:short-circuits-safety:3}
De omleiding steelt de hele mars van \omterm{def:g3:first-electric-circuit:bulb}{lampje} 2 --- in de steek
gelaten gaat het uit. Met één tegenwerker overgeslagen werkt de hele
lus de \omterm{def:g3:first-electric-circuit:battery}{batterij} minder tegen, versterkt de mars, en wordt
\omterm{def:g3:first-electric-circuit:bulb}{lampje} 1 --- dat de sterkere stroom nu alleen draagt --- feller.
\end{solution}

\begin{solution}{exo:g7:short-circuits-safety:4}
Een marcherende stroom verwarmt zijn weg altijd, en een stormloop
verwarmt hem razend. De lamp gebruikt die verwarming met opzet: haar
draadje is gebouwd om witheet te lopen en te schijnen.
\end{solution}

\begin{solution}{exo:g7:short-circuits-safety:5}
Een met opzet dun opofferingsdraadje dat \omterm{def:g2:ice-water-steam:meltfreeze}{smelt} --- en zo de lus
doorsnijdt --- zodra de stroom zijn waarde overschrijdt; het offert
zichzelf op. \omterm{def:g4:series-circuits:series}{In serie}, omdat de stormloop door de controlepost
gedwongen moet worden: een wachter \omterm{def:g5:series-parallel-circuits:parallel}{parallel} zou hem alleen zien
voorbijgaan.
\end{solution}

\begin{solution}{exo:g7:short-circuits-safety:6}
De onderbreker klapt open in plaats van te \omterm{def:g6:water-states:changes}{smelten}, en wordt met een
tik teruggezet --- geen reserveonderdelen. Vóór elk terugzetten of
vervangen moet de storing gevonden en verholpen zijn: de wachter is
niet zomaar uit gevlogen.
\end{solution}

\begin{solution}{exo:g7:short-circuits-safety:7}
Een overbelasting: drie onschuldige apparaten vroegen één
\omterm{def:g5:series-parallel-circuits:parallel}{tak} om hun drie volle marsen tegelijk te dragen, en de opgetelde
stormloop verwarmt de bedrading. De onderbreker klapt op de som eruit
--- terecht.
\end{solution}

\begin{solution}{exo:g7:short-circuits-safety:8}
Een spijker \omterm{def:g2:ice-water-steam:meltfreeze}{smelt} nooit: hij legt de wachter het zwijgen op terwijl
de ziekte woedt. De volgende stormloop treft geen controlepost aan, en
de wandbedrading van het huis zelf wordt de \omterm{def:g7:short-circuits-safety:fuse}{zekering} --- verkolend
waar niemand het kan zien.
\end{solution}

\begin{solution}{exo:g7:short-circuits-safety:9}
Een \omterm{def:g7:short-circuits-safety:device}{kortsluiting} stroomopwaarts van het apparaat --- feitelijk een
\omterm{def:g7:short-circuits-safety:device}{kortsluiting} over de voeding, gezien vanaf dat snoer. De lamp sterft
(in de steek gelaten), de stormloop raast door het gerafelde contact
en verhit de kabel --- en de wachter van de \omterm{def:g5:series-parallel-circuits:parallel}{tak},
\omterm{def:g7:short-circuits-safety:fuse}{zekering} of onderbreker, snijdt de \omterm{def:g5:series-parallel-circuits:parallel}{tak} af. Die snee is het
systeem dat werkt.
\end{solution}

\begin{solution}{exo:g7:short-circuits-safety:10}
Over allebei de \omterm{def:g3:first-electric-circuit:bulb}{lampjes} gelegd verbindt de draad --- via gewone draad
en de aansluitingen van de \omterm{def:g3:first-electric-circuit:battery}{batterij} --- de eigen twee
polen van de \omterm{def:g3:first-electric-circuit:battery}{batterij} zonder enig apparaat op de omleidingsweg:
een \omterm{def:g7:short-circuits-safety:device}{kortsluiting} van de \omterm{def:g3:first-electric-circuit:battery}{batterij}. Beide \omterm{def:g3:first-electric-circuit:bulb}{lampjes} sterven en
de stormloop verhit draad en \omterm{def:g3:first-electric-circuit:battery}{batterij}; zelfs op speelgoedschaal
verwoest dat cellen en verwarmt het vingers, dus wordt het in woorden
gedemonstreerd en niet in daden.
\end{solution}

\begin{solution}{exo:g7:short-circuits-safety:11}
De aan elkaar gesmolten steundraadjes van het stervende
\omterm{def:g3:first-electric-circuit:bulb}{lampje} werden een ingebouwde omleiding: de lus blijft gesloten
\emph{door} de \omterm{def:g7:short-circuits-safety:device}{kortsluiting} van het dode \omterm{def:g3:first-electric-circuit:bulb}{lampje}, dus overleeft de
sliert zijn verlies --- één tegenwerker minder, dus iets feller. Maar
elke zo'n dood versterkt de mars door de overlevenden en haast de
volgende uitval: de sliert bereidt beleefd haar eigen kettingreactie
voor.
\end{solution}

\begin{solution}{exo:g7:short-circuits-safety:12}
Eén hoofdwachter en hoofdschakelaar op de gemeenschappelijke weg; dan
elke \omterm{def:g5:series-parallel-circuits:parallel}{tak} --- verlichting, gereedschap, lader --- met een eigen
\omterm{ex:g3:first-electric-circuit:switch}{schakelaar} en een eigen wachter op maat van zijn \omterm{def:g5:series-parallel-circuits:parallel}{tak} (de
\omterm{def:g5:series-parallel-circuits:parallel}{tak} van het gereedschap draagt de zwaarste marsen en verdient de
stevigste onderbreker; de dunne \omterm{def:g5:series-parallel-circuits:parallel}{tak} van de lader wil de meest
tere \omterm{def:g7:short-circuits-safety:fuse}{zekering}, want die zou verkolen bij stromen waar de hoofdwachter
zich niets van \omterm{def:g1:magnets:magnet}{aantrekt}). Bordje boven de lade, bijvoorbeeld: ``Een
\omterm{def:g7:short-circuits-safety:fuse}{zekering} die springt vertelt de waarheid --- zoek de storing. Vervang
gelijk door gelijk; folie en spijkers bewaken niets dan de weg van het
vuur.''
\end{solution}

\begin{solution}{pb:g7:short-circuits-safety:1}
\textbf{1.} Een overbelasting: drie warmtehongerige apparaten telden
hun eerlijke marsen op één \omterm{def:g5:series-parallel-circuits:parallel}{tak} op --- er was geen blote draad nodig;
de som alleen al verwarmt bedrading tot voorbij de veiligheid.
\textbf{2.} Een \omterm{def:g7:short-circuits-safety:device}{kortsluiting} bij de beet: vanaf dat moment liet de
mars de lamp erachter in de steek en stormde hij door de elkaar
rakende \omterm{def:g4:conductors-insulators:conductor}{geleiders}, waarbij het snoer verhitte.
\textbf{3.} Ze schafte haar af: folie geleidt elke stormloop zonder te
\omterm{def:g6:water-states:changes}{smelten} --- de \omterm{def:g5:series-parallel-circuits:parallel}{tak} stond zonder wachter.
\textbf{4.} Waarschijnlijk verhaal: de overbelasting kwam eerst, en de
echte \omterm{def:g7:short-circuits-safety:fuse}{zekering} sprong --- eerlijk --- misschien meer dan eens; iemand
``genas'' het ongemak met folie (stuk drie); toen de \omterm{def:g7:short-circuits-safety:device}{kortsluiting} van
het aangevreten snoer eindelijk toesloeg, was er geen wachter meer
over, en verhitte de stormloop de wandbedrading tot brand in plaats
van te sterven bij een gesprongen \omterm{def:g7:short-circuits-safety:fuse}{zekering} in het donker.
\textbf{5.} Een stroom verwarmt zijn weg altijd; harder stormen,
feller verwarmen. De lamp bewijst dat verwarming te temmen valt ---
haar draadje is ontworpen voor witte gloed in een veilig
\omterm{def:g3:first-electric-circuit:bulb}{lampje}. Wat hier ontbrak: elke ontworpen plek waar de
\omterm{def:g5:heat-and-insulation:heat}{warmte} mocht wonen, en elke wachter om haar te stoppen.
\textbf{6.} Met opzet dun, zodat juist zij van alle draden van de
kring als eerste \omterm{def:g2:ice-water-steam:meltfreeze}{smelt}: de wachter neemt de fatale verhitting in
haar eigen drie \omterm{def:g2:measuring-length:units}{centimeter} --- ze sterft opdat de \omterm{def:g2:measuring-length:units}{meters} in de muren
dat niet doen.
\textbf{7.} Een wachter moet op het juiste moment \emph{falen}: haar
deugd is een ontworpen zwakte. De voortreffelijkheid van folie als
\omterm{def:g4:conductors-insulators:conductor}{geleider} was precies de corruptie van de controlepost --- perfecte
doorgang, geen bescherming.
\textbf{8.} Veiligheid telt niet op als beleefdheid: drie rechtmatige
marsen op één weg tellen op tot één onrechtmatige stormloop --- de
\omterm{def:g5:series-parallel-circuits:parallel}{tak} gehoorzaamt aan rekenkunde, niet aan bedoelingen.
\textbf{9.} De folie. Overbelastingen en aangevreten snoeren zijn
ongelukken die het systeem geacht wordt te overleven --- wachters
bestaan juist daarvoor; de wachter het zwijgen opleggen maakte van
overleefbare storingen een ramp.
\textbf{10.} Met een deugdelijke \omterm{def:g7:short-circuits-safety:fuse}{zekering}: de \omterm{def:g7:short-circuits-safety:device}{kortsluiting} van het
snoer slaat toe, de stroom springt op, de \omterm{def:g7:short-circuits-safety:fuse}{zekering} \omterm{def:g2:ice-water-steam:meltfreeze}{smelt} binnen
ogenblikken --- de keuken wordt donker, het gezin komt thuis bij een
gesprongen \omterm{def:g7:short-circuits-safety:fuse}{zekering} en een aangevreten snoer, moppert, repareert
allebei, en komt er nooit achter hoe erg de nacht bijna was.
\textbf{11.} Elk warmteapparaat aan zijn eigen wandcontact op zijn
eigen afdoende bewaakte \omterm{def:g5:series-parallel-circuits:parallel}{tak}; de stekkerdoos terug naar alleen lampen
en opladers, onder de vaste regel: één hongerig apparaat per
stopcontact, en nooit voedt een stekkerdoos een kachel.
\textbf{12.} Bijvoorbeeld: ``De werkwijze van de schurk is de
moeiteloze weg: bied de mars een omleiding en hij neemt haar, en
verhit terwijl hij stormt. De taak van de wachter is \omterm{def:g4:series-circuits:series}{in serie} te
staan en als eerste te falen, eerlijk en goedkoop. En geen huishouden
mag ooit een wachter improviseren: een \omterm{def:g7:short-circuits-safety:fuse}{zekering} wordt vervangen door
een \omterm{def:g7:short-circuits-safety:fuse}{zekering} --- gelijk voor gelijk --- of de muren zelf zullen op een
dag het \omterm{def:g6:water-states:changes}{smelten} voor hun rekening nemen.''
\end{solution}
